\begin{proofbox}
	\textbf{\textcolor{CProof}{思路提示}}

	通过圆心角作为桥梁：等弧 $\Leftrightarrow$ 等圆心角 $\Leftrightarrow$ 等圆周角。

	\medskip
	\textbf{\textcolor{CProof}{正式推导}}
	\begin{description}[style=nextline, leftmargin=2.8em, labelsep=0.5em, itemsep=0.55em]
		\item[\textbf{步骤一（弧等推圆心角等）：}]
			\[
				\overset{\frown}{AB} = \overset{\frown}{CD}.
			\]
			\[
				\angle AOB = \angle COD.
			\]
			\par\textbf{理由：} 在同圆或等圆中，等弧所对的圆心角相等。

		\item[\textbf{步骤二（圆心角转圆周角）：}]
			\[
				\angle APB = \frac{1}{2}\angle AOB.
			\]
			\[
				\angle CQD = \frac{1}{2}\angle COD.
			\]
			\par\textbf{理由：} 同弧所对的圆周角等于圆心角的一半。

		\item[\textbf{步骤三（等量代换）：}]
			\[
				\angle APB = \angle CQD.
			\]
			\par\textbf{理由：} 将步骤二的关系代入步骤一的等式中。

		\item[\textbf{步骤四（圆周角等推圆心角等）：}]
			\[
				\angle AOB = 2\angle APB,\quad \angle COD = 2\angle CQD.
			\]
			\par\textbf{理由：} 由步骤二变形。因为$\angle APB = \angle CQD$，所以$\angle AOB = \angle COD$。

		\item[\textbf{步骤五（圆心角等推弧等）：}]
			\[
				\angle AOB = \angle COD.
			\]
			\[
				\overset{\frown}{AB} = \overset{\frown}{CD}.
			\]
			\par\textbf{理由：} 在同圆或等圆中，相等的圆心角所对的弧相等。

		\item[\textbf{步骤六（归纳双向）：}]
			\textbf{充分性：} 若$\overset{\frown}{AB} = \overset{\frown}{CD}$，则$\angle APB = \angle CQD$；
			\textbf{必要性：} 若$\angle APB = \angle CQD$，则$\overset{\frown}{AB} = \overset{\frown}{CD}$。
			\par\textbf{理由：} 充分性由步骤一至三构建，必要性由步骤四至五构建。
	\end{description}

	\medskip
	\textbf{\textcolor{CProof}{结论回扣}}

	\textbf{上述推导表明，在同圆或等圆中，弧相等与圆周角相等互为充要条件，可用于角与弧的相互推导。}
\end{proofbox}
